\documentclass{ltugboat}
\usepackage{amsmath}
\usepackage{microtype}
\usepackage[hidelinks]{hyperref}
\def\url{\tburl}
\title{Reviewing an assisted edit to an existing \LaTeX\ project}
% Editorial draft: author attribution and correspondence to be confirmed
% by the invited author before submission. No manuscript has been sent.
\author{Fermion development team}
\address{Fermion Inc., Japan}
\netaddress{contact (at) fermion dot company}
\begin{document}
\maketitle
\begin{abstract}
An assisted edit should leave enough evidence to distinguish a change to
source from a successful build and from a correct mathematical argument.
We describe a small, reproducible workflow using an existing two-file
LaTeX project. A centered-difference calculation supplies a deliberate
error that compiles successfully. Source differences, an independent
polynomial calculation, and inspection of the resulting PDF provide
separate checks. The example also records an interrupted assistant turn
and the distinction between a PDF on disk and the PDF currently shown
in the editor.
\end{abstract}

\section{The review boundary}
A mathematically wrong expression can be perfectly valid \LaTeX.
Conversely, a correct edit can fail to build because a file or package is
missing. An editor can also retain a previously successful \acro{PDF}
after a failed or interrupted run. These cases call for different
checks; a single green status or fluent explanation cannot distinguish
them.

We use TeX64, a desktop editor for \macOS{} and Windows, as a concrete
example. Its Axiom assistant can read and edit files in the opened
project, invoke a build, and inspect a \acro{PDF}. This article concerns
the evidence left by those actions rather than a comparison of models
or editing speed. The original fixture, the exact prompt, the actual
edited source, and the source difference accompany the article.
The argument and the build can be checked without access to Axiom.

The authoring tool and the reviewing tool need not be the same.
We keep the ordinary \texttt{.tex} files as the authoritative manuscript.
The assistant's explanation is an observation about that manuscript,
not a replacement for it. In particular, a statement that an equation
was checked does not establish which revision was checked.

\section{A deliberately wrong but compilable example}
The fixture has this structure:
\begin{verbatim}[\small]
project/
  main.tex
  sections/stencil.tex
\end{verbatim}
The main file loads \texttt{amsmath}, sets a title, and uses
\verb|\input{sections/stencil}|. It requires no external images,
bibliography, shell escape, or nonstandard document class. It is an
original fictional manuscript, not a student's answer or a client's
paper. The section file defines
\begin{equation}\label{eq:dh}
 D_hu(x)=\frac{u(x+h)-2u(x)+u(x-h)}{h^2}.
\end{equation}
For a function with six continuous derivatives near $x$, the fixture
then incorrectly prints
\begin{equation}\label{eq:bad}
 D_hu(x)=u''(x)+\frac{h}{12}u^{(4)}(x)+O(h^4).
\end{equation}
The prose says that the leading error is second order. Both the wrong
formula and the inconsistent prose compile. The \TeX{} log therefore
cannot serve as a mathematical correctness check.

Taylor expansion gives
\begin{align*}
 u(x\mathbin{\pm}h)
 &=u(x)\mathbin{\pm}hu'(x)+\frac{h^2}{2}u''(x)\\
 &\quad\mathbin{\pm}\frac{h^3}{6}u^{(3)}(x)
 +\frac{h^4}{24}u^{(4)}(x)\\
 &\quad\mathbin{\pm}\frac{h^5}{120}u^{(5)}(x)+O(h^6).
\end{align*}
The sum cancels the odd powers. After subtracting $2u(x)$ and dividing
by $h^2$, the corrected expression is
\begin{equation}\label{eq:good}
 D_hu(x)=u''(x)+\frac{h^2}{12}u^{(4)}(x)+O(h^4).
\end{equation}
This exposes the local edit clearly: the numerator of the leading
error term must contain $h^2$, not $h$.

\section{An independent numerical check}
Take $u(x)=x^4$, $x=1$, and $h=1/2$. Exact rational arithmetic gives
\begin{align*}
 D_{1/2}u(1)
 &=\frac{(3/2)^4-2+(1/2)^4}{(1/2)^2}
 =\frac{25}{2},\\
 u''(1)&=12,\\
 D_{1/2}u(1)-u''(1)&=\frac12.
\end{align*}
Since $u^{(4)}(1)=24$, the corrected leading term is
$(1/2)^2\,24/12=1/2$. The wrong term would give $1$.
Here the higher-order terms vanish because the function is a quartic.
This check is intentionally small and exact. It tests the disputed
power of $h$ without introducing floating-point tolerances.

A second rational value, such as $h=1/4$, gives an error of $1/8$.
Halving $h$ divides the error by four for this polynomial. That fact
illustrates second-order behavior for the fixture; it is not a timing
measurement or a general empirical accuracy benchmark.

\section{The assisted editing run}
We opened the existing fixture folder in the installed TeX64 0.1.26
on \macOS. The original source was first built and its \acro{PDF}
retained. The request to Axiom specified the section file, the equation
to repair, the independent polynomial check, and the limits of the
edit. Its operative instructions were to derive the Taylor expansion,
correct the leading power of $h$, preserve labels and project
structure, build \texttt{main.tex}, and inspect the resulting
\acro{PDF}. The exact request is in \texttt{axiom-prompt.txt}.

In the section file, Axiom added the two Taylor expansions and the
exact error calculation, replacing $h/12$ by
$h^2/12$. The actual difference reports 23 added lines and 8 removed
lines. The main document, packages, title, and existing equation labels
were preserved. A separately written reference solution is supplied;
This reference was prepared independently of the observed edit and is
identified separately from the assistant output.

The observed turn stopped at its processing limit before another
model call. The application exposed a continuation action. A corrected
\acro{PDF} existed on disk, while an earlier \acro{PDF} remained visible
in the editor during inspection. We therefore did not infer successful
visual review from the presence of the edited file or the build output.
The continuation again displayed a processing-limit message, but the
final editor view showed the corrected formula and polynomial check.
The supplementary record distinguishes these observations. An interruption
is useful evidence about the workflow, even when the source edit itself
is correct.

\section{Retaining reviewable artifacts}
Before editing, copy the project to a new directory. Preserve the
original source and original \acro{PDF}; do not use an old output as
proof of a new build. After the run, retain these artifacts:
\begin{enumerate}
\item The exact original inputs and assistant request.
\item The actual edited inputs and a unified source difference.
\item The command, compiler version, exit status, and build log.
\item The new \acro{PDF}, its checksum, and the pages inspected.
\item The independent mathematical calculation and its assumptions.
\end{enumerate}
These items answer different questions. The difference shows what
changed. The log shows whether the chosen inputs compiled. The
checksum identifies which output was inspected. A visual check can
find clipped equations or stale references. The polynomial calculation
checks a mathematical consequence of the change.

A byte-identical \acro{PDF} is not required across all machines:
metadata, fonts, and compiler versions can differ. Checksums identify
local artifacts, not portable mathematical equivalence. Retaining the
source and compiler information makes those differences investigable.
A successful build also does not establish that an arbitrary document
is correct; this example has an explicit, independently checkable
acceptance condition.

\section{Building without the assistant}
The supplementary fixture can be built with an ordinary \TeX{}
distribution. In a copy of the original project, replace the section
file by \path{axiom-stencil.tex}, the actual observed edit. Then run
these commands from the directory containing \texttt{main.tex}:
\begin{verbatim}[\small]
pdflatex -interaction=nonstopmode \
  -halt-on-error main.tex
pdflatex -interaction=nonstopmode \
  -halt-on-error main.tex
\end{verbatim}
The second pass resolves the equation references. Do not enable shell
escape for this fixture. Check the exit codes and read the log for
undefined references, missing files, and overfull boxes. Open the
resulting \acro{PDF} and compare equations (1) and (2), the exponent in
the error term, and the final value $1/2$ with the source.

The source comparison can be produced independently of the editor:
\begin{verbatim}[\small]
diff -u before/stencil.tex \
  after/stencil.tex
shasum -a 256 main.pdf
\end{verbatim}
The supplementary record contains the commands actually run in our
environment. The commands above describe the portable reproduction;
they do not imply that a Windows run was performed. Windows behavior
was not tested in this preparation.

\section{Continuing an exported project}
An Overleaf project can be downloaded as a source archive and unpacked
before opening its folder in a local editor. Preserve relative file
paths, the main document selection, the compiler choice, images, and
bibliography inputs. A source archive is not the same as the project's
compiled \acro{PDF}. Missing fonts or packages may need to be supplied
locally. Overleaf's documentation describes its source-download
procedure~\cite{overleaf}.

For a migration check, first build the unchanged export. Then perform
one bounded edit and compare both the source difference and the new
\acro{PDF}. This separates an environment mismatch from an editing
error. It also provides a usable stopping point if an assisted turn is
interrupted.

Local editing and compilation differ from the assistant's network
operation. Axiom sends document context needed for its request to the
service. This workflow therefore uses a fictional manuscript. It
should not be described as entirely offline. Network access is also
needed for an initial managed \TeX{} download; an existing installed
\TeX{} distribution can be used instead.

\section{Scope and limitations}
This is one small, deliberately constructed example. It establishes
that the captured source change can be reviewed, compiled, and checked
against exact arithmetic. It does not estimate general assistant
accuracy, saved time, energy use, adoption, or comparative performance.
The interrupted turn and stale displayed output must remain in the
record; removing them would hide an actual failure mode.

The authoring team develops TeX64. The article and fixture are intended
to explain a reproducible review procedure and disclose its limits.
There are no purchase links or claims that a subscription is needed to
reproduce the mathematical result. Larger projects, bibliographies,
custom fonts, external build tools, and platform-specific behavior need
their own recorded checks.

\begin{thebibliography}{9}
\bibitem{overleaf} Overleaf.
Downloading a project.\newline
\url{https://docs.overleaf.com/managing-projects-and-files/downloading-a-project}.
\bibitem{tex64} Fermion Inc. TeX64 documentation.\newline
\url{https://tex64.com/docs/}.
\bibitem{tugboat} \TeX{} Users Group.
\TUB{} author information and macros.\newline
\url{https://tug.org/tugboat/location.html}.
\end{thebibliography}
\makesignature
\end{document}
